% GENERATED FILE. Edit canonical lessons, then run npm run generate:books.
% canonical-source-hash: fnv1a64:21f8745ceacc3eb5
% canonical-lessons: HI-C300-dhadhakna, HI-C300-dahakna, HI-C300-jvala, HI-C300-ujala, HI-C300-tapis

\chapter{The Blaze}
\label{ch:hi-the-blaze}
\hlchaptermodality{300}

\begin{chapteropening}
\textbf{By the end of this chapter:} \emph{I can describe a fire at its height.}

Complete the last lesson of chapter 300: I can describe a fire at its height.
\end{chapteropening}

\section[\texorpdfstring{\hi{धधकना} (dhadhaknā)}{धधकना (dhadhaknā)}]{\hi{धधकना} (dhadhaknā) — to blaze}

\label{lesson:HI-C300-dhadhakna}



\begin{quote}
Before the new one: say the Hindi for to flare up, then the Hindi for dry straw.
\end{quote}

\subsection*{You'll want to know: \hi{धधकना}}
\textbf{\hi{धधकना}} — \emph{dhadhaknā} — \textquotedblleft{}to blaze\textquotedblright{}.

\begin{quote}
\hi{पूरा} \hi{घर} \hi{धधक} \hi{रहा} \hi{था।}
\end{quote}

\begin{quote}
\emph{pūrā ghar dhadhak rahā thā.}
\end{quote}

\begin{quote}
The whole house was ablaze.
\end{quote}

Past (it …): \hi{वह} \hi{आग} \hi{धधकी}, \emph{vah āg dhadhkī}.

\begin{cousinweb}
\textbf{\hi{धधकना}} (\emph{dhadhaknā}) is inherited from \emph{*dhagg}, a form CDIAL reconstructs (CDIAL 6704); the doubled \emph{dha-dha} gives the roar of a big fire.
\end{cousinweb}

\begin{grammarlens}[title={What you've built}]
One more word for telling what happened.
\end{grammarlens}

\subsection*{Your turn}
\begin{itemize}
\raggedright
  \item \emph{Say it:} \emph{dhadhaknā}
  \item \emph{Say it:} \emph{dhadhaknā}, once more
  \item \emph{Say it:} \emph{pūrā ghar dhadhak rahā thā.} — The whole house was ablaze.
  \item \emph{Recall:} say the Hindi for to flare up, then the Hindi for dry straw, then say \emph{dhadhaknā} again
\end{itemize}

\begin{tcolorbox}[breakable,colback=teal!4,colframe=teal!35!black,title={Before you move on}]
What does \hi{धधकना} mean? (\textquotedblleft{}To blaze\textquotedblright{}.) Say it once more.
\end{tcolorbox}
\section[\texorpdfstring{\hi{दहकना} (dahaknā)}{दहकना (dahaknā)}]{\hi{दहकना} (dahaknā) — to glow red-hot}

\label{lesson:HI-C300-dahakna}



\begin{quote}
Before the new one: say the Hindi for dry straw, then the Hindi for to blaze.
\end{quote}

\subsection*{You'll want to know: \hi{दहकना}}
\textbf{\hi{दहकना}} — \emph{dahaknā} — \textquotedblleft{}to glow red-hot\textquotedblright{}.

\begin{quote}
\hi{आग} \hi{के} \hi{बाद} \hi{भी} \hi{अंगारे} \hi{दहक} \hi{रहे} \hi{थे।}
\end{quote}

\begin{quote}
\emph{āg ke bād bhī aṅgāre dahak rahe the.}
\end{quote}

\begin{quote}
Even after the fire the embers were still glowing.
\end{quote}

Past (it …): \hi{वह} \hi{कोयला} \hi{दहका}, \emph{vah koylā dahkā}.

\begin{cousinweb}
\textbf{\hi{दहकना}} (\emph{dahaknā}) goes back to Sanskrit \emph{dahati}, \textquotedblleft{}burns\textquotedblright{} (CDIAL 6245). It is the glow of coals with no flame: embers \textbf{\hi{दहकते} \hi{हैं}}.
\end{cousinweb}

\begin{grammarlens}[title={What you've built}]
One more word for telling what happened.
\end{grammarlens}

\subsection*{Your turn}
\begin{itemize}
\raggedright
  \item \emph{Say it:} \emph{dahaknā}
  \item \emph{Say it:} \emph{dahaknā}, once more
  \item \emph{Say it:} \emph{āg ke bād bhī aṅgāre dahak rahe the.} — Even after the fire the embers were still glowing.
  \item \emph{Recall:} say the Hindi for dry straw, then the Hindi for to blaze, then say \emph{dahaknā} again
\end{itemize}

\begin{tcolorbox}[breakable,colback=teal!4,colframe=teal!35!black,title={Before you move on}]
What does \hi{दहकना} mean? (\textquotedblleft{}To glow red-hot\textquotedblright{}.) Say it once more.
\end{tcolorbox}
\section[\texorpdfstring{\hi{ज्वाला} (jvālā)}{ज्वाला (jvālā)}]{\hi{ज्वाला} (jvālā) — a blaze; a flame}

\label{lesson:HI-C300-jvala}



\begin{quote}
Before the new one: say the Hindi for to blaze, then the Hindi for to glow red-hot.
\end{quote}

\subsection*{You'll want to know: \hi{ज्वाला}}
\textbf{\hi{ज्वाला}} — \emph{jvālā} — \textquotedblleft{}a blaze; a flame\textquotedblright{}.

\begin{quote}
\hi{आग} \hi{की} \hi{ज्वाला} \hi{दूर} \hi{से} \hi{दिख} \hi{रही} \hi{थी।}
\end{quote}

\begin{quote}
\emph{āg kī jvālā dūr se dikh rahī thī.}
\end{quote}

\begin{quote}
The blaze could be seen from far away.
\end{quote}

Feminine.

\begin{cousinweb}
\textbf{\hi{ज्वाला}} (\emph{jvālā}) is Sanskrit \emph{jvālā}, \textquotedblleft{}a flame\textquotedblright{} (Monier-Williams), taken over whole. It is a word for writing and for big fires; the everyday word is \textbf{\hi{लपट}}.
\end{cousinweb}

\begin{grammarlens}[title={What you've built}]
One more word for telling what happened.
\end{grammarlens}

\subsection*{Your turn}
\begin{itemize}
\raggedright
  \item \emph{Say it:} \emph{jvālā}
  \item \emph{Say it:} \emph{jvālā}, once more
  \item \emph{Say it:} \emph{āg kī jvālā dūr se dikh rahī thī.} — The blaze could be seen from far away.
  \item \emph{Recall:} say the Hindi for to blaze, then the Hindi for to glow red-hot, then say \emph{jvālā} again
\end{itemize}

\begin{tcolorbox}[breakable,colback=teal!4,colframe=teal!35!black,title={Before you move on}]
What does \hi{ज्वाला} mean? (\textquotedblleft{}A blaze; a flame\textquotedblright{}.) Say it once more.
\end{tcolorbox}
\section[\texorpdfstring{\hi{उजाला} (ujālā)}{उजाला (ujālā)}]{\hi{उजाला} (ujālā) — light; brightness}

\label{lesson:HI-C300-ujala}



\begin{quote}
Before the new one: say the Hindi for to glow red-hot, then the Hindi for a blaze.
\end{quote}

\subsection*{You'll want to know: \hi{उजाला}}
\textbf{\hi{उजाला}} — \emph{ujālā} — \textquotedblleft{}light; brightness\textquotedblright{}.

\begin{quote}
\hi{आग} \hi{का} \hi{उजाला} \hi{पूरे} \hi{गाँव} \hi{में} \hi{फैल} \hi{गया।}
\end{quote}

\begin{quote}
\emph{āg kā ujālā pūre gā̃v meṁ phail gayā.}
\end{quote}

\begin{quote}
The glow of the fire spread over the whole village.
\end{quote}

Masculine.

\begin{cousinweb}
\textbf{\hi{उजाला}} (\emph{ujālā}) is inherited from \emph{*ujjvālaka}, \textquotedblleft{}bright\textquotedblright{}, a form CDIAL reconstructs (CDIAL 1673). It is light that lets you see, the opposite of darkness: \textbf{\hi{सुबह} \hi{का} \hi{उजाला}}, the light of morning.
\end{cousinweb}

\begin{grammarlens}[title={What you've built}]
One more word for telling what happened.
\end{grammarlens}

\subsection*{Your turn}
\begin{itemize}
\raggedright
  \item \emph{Say it:} \emph{ujālā}
  \item \emph{Say it:} \emph{ujālā}, once more
  \item \emph{Say it:} \emph{āg kā ujālā pūre gā̃v meṁ phail gayā.} — The glow of the fire spread over the whole village.
  \item \emph{Recall:} say the Hindi for to glow red-hot, then the Hindi for a blaze, then say \emph{ujālā} again
\end{itemize}

\begin{tcolorbox}[breakable,colback=teal!4,colframe=teal!35!black,title={Before you move on}]
What does \hi{उजाला} mean? (\textquotedblleft{}Light; brightness\textquotedblright{}.) Say it once more.
\end{tcolorbox}
\section[\texorpdfstring{\hi{तपिश} (tapiś)}{तपिश (tapiś)}]{\hi{तपिश} (tapiś) — scorching heat}

\label{lesson:HI-C300-tapis}



\begin{quote}
Before the new one: say the Hindi for a blaze, then the Hindi for light.
\end{quote}

\subsection*{You'll want to know: \hi{तपिश}}
\textbf{\hi{तपिश}} — \emph{tapiś} — \textquotedblleft{}scorching heat\textquotedblright{}.

\begin{quote}
\hi{आग} \hi{की} \hi{तपिश} \hi{से} \hi{खिड़की} \hi{का} \hi{शीशा} \hi{टूट} \hi{गया।}
\end{quote}

\begin{quote}
\emph{āg kī tapiś se khiṛkī kā śīśā ṭūṭ gayā.}
\end{quote}

\begin{quote}
The heat of the fire cracked the window glass.
\end{quote}

Feminine.

\begin{cousinweb}
\textbf{\hi{तपिश}} (\emph{tapiś}) is usually traced to Persian \emph{tapiš}, \textquotedblleft{}heat, a burning\textquotedblright{}. It is stronger than \textbf{\hi{गर्मी}}: the heat that comes off a fire or a road at noon.
\end{cousinweb}

\begin{grammarlens}[title={What you've built}]
One more word for telling what happened.
\end{grammarlens}

\subsection*{Your turn}
\begin{itemize}
\raggedright
  \item \emph{Say it:} \emph{tapiś}
  \item \emph{Say it:} \emph{tapiś}, once more
  \item \emph{Say it:} \emph{āg kī tapiś se khiṛkī kā śīśā ṭūṭ gayā.} — The heat of the fire cracked the window glass.
  \item \emph{Recall:} say the Hindi for a blaze, then the Hindi for light, then say \emph{tapiś} again
\end{itemize}

\begin{tcolorbox}[breakable,colback=teal!4,colframe=teal!35!black,title={Before you move on}]
What does \hi{तपिश} mean? (\textquotedblleft{}Scorching heat\textquotedblright{}.) Say it once more.
\end{tcolorbox}
